% GENERATED FILE. Edit canonical lessons, then run npm run generate:books.
% canonical-source-hash: fnv1a64:486575b5a00dff42
% canonical-lessons: SA-C378-lesah, SA-C378-manak, SA-C378-ardham, SA-C378-nyunah, SA-C378-hrasati

\chapter{A Little}
\label{ch:sa-a-little}
\hlchaptermodality{378}

\begin{chapteropening}
\textbf{By the end of this chapter:} \emph{I can say what a small change would do.}

Complete the last lesson of chapter 378: I can say what a small change would do.
\end{chapteropening}

\section[\texorpdfstring{\sk{लेशः} (leśaḥ)}{लेशः (leśaḥ)}]{\sk{लेशः} (leśaḥ) — a tiny portion; a little bit}

\label{lesson:SA-C378-lesah}



\begin{quote}
Before the new one: say the Sanskrit for completely, then the Sanskrit for exceeds.
\end{quote}

\subsection*{You'll want to know: \sk{लेशः}}
\textbf{\sk{लेशः}} — \emph{leśaḥ} — \textquotedblleft{}a tiny portion; a little bit\textquotedblright{}.

\begin{quote}
\sk{यदि} \sk{दोषस्य} \sk{लेशः} \sk{अपि} \sk{अस्ति}, \sk{तर्हि} \sk{पुनः} \sk{चिन्तयतु।}
\end{quote}

\begin{quote}
\emph{yadi doṣasya leśaḥ api asti, tarhi punaḥ cintayatu.}
\end{quote}

\begin{quote}
If there is even a little of a mistake in it, think again.
\end{quote}

A masculine noun.

\begin{cousinweb}
\textbf{\sk{लेशः}} (\emph{leśaḥ}) MW gives \emph{leśa} \textquotedblleft{}a small part or portion, a particle, a little bit\textquotedblright{} and gives no derivation, so its origin is uncertain (MW).
\end{cousinweb}

\begin{grammarlens}[title={What you've built}]
One more word for saying what happens if.
\end{grammarlens}

\subsection*{Your turn}
\begin{itemize}
\raggedright
  \item \emph{Say it:} \emph{leśaḥ}
  \item \emph{Say it:} \emph{leśaḥ}, once more
  \item \emph{Say it:} \emph{yadi doṣasya leśaḥ api asti, tarhi punaḥ cintayatu.} — If there is even a little of a mistake in it, think again.
  \item \emph{Recall:} say the Sanskrit for completely, then the Sanskrit for exceeds, then say \emph{leśaḥ} again
\end{itemize}

\begin{tcolorbox}[breakable,colback=teal!4,colframe=teal!35!black,title={Before you move on}]
What does \sk{लेशः} mean? (\textquotedblleft{}A tiny portion; a little bit\textquotedblright{}.) Say it once more.
\end{tcolorbox}
\section[\texorpdfstring{\sk{मनाक्} (manāk)}{मनाक् (manāk)}]{\sk{मनाक्} (manāk) — a little; in a small degree}

\label{lesson:SA-C378-manak}



\begin{quote}
Before the new one: say the Sanskrit for exceeds, then the Sanskrit for a tiny portion.
\end{quote}

\subsection*{You'll want to know: \sk{मनाक्}}
\textbf{\sk{मनाक्}} — \emph{manāk} — \textquotedblleft{}a little; in a small degree\textquotedblright{}.

\begin{quote}
\sk{यदि} \sk{सः} \sk{मनाक्} \sk{अपि} \sk{यतते}, \sk{तर्हि} \sk{सफलतां} \sk{प्राप्नोति।}
\end{quote}

\begin{quote}
\emph{yadi saḥ manāk api yatate, tarhi saphalatāṃ prāpnoti.}
\end{quote}

\begin{quote}
If he tries even a little, he succeeds.
\end{quote}

\begin{cousinweb}
\textbf{\sk{मनाक्}} (\emph{manāk}) MW says \emph{manāk} probably comes from \emph{manā} and the root √añc (\textquotedblleft{}perceivably\textquotedblright{}), and gives \textquotedblleft{}a little, slightly, in a small degree\textquotedblright{} (MW).
\end{cousinweb}

\begin{grammarlens}[title={What you've built}]
One more word for saying what happens if.
\end{grammarlens}

\subsection*{Your turn}
\begin{itemize}
\raggedright
  \item \emph{Say it:} \emph{manāk}
  \item \emph{Say it:} \emph{manāk}, once more
  \item \emph{Say it:} \emph{yadi saḥ manāk api yatate, tarhi saphalatāṃ prāpnoti.} — If he tries even a little, he succeeds.
  \item \emph{Recall:} say the Sanskrit for exceeds, then the Sanskrit for a tiny portion, then say \emph{manāk} again
\end{itemize}

\begin{tcolorbox}[breakable,colback=teal!4,colframe=teal!35!black,title={Before you move on}]
What does \sk{मनाक्} mean? (\textquotedblleft{}A little; in a small degree\textquotedblright{}.) Say it once more.
\end{tcolorbox}
\section[\texorpdfstring{\sk{अर्धम्} (ardham)}{अर्धम् (ardham)}]{\sk{अर्धम्} (ardham) — a half}

\label{lesson:SA-C378-ardham}



\begin{quote}
Before the new one: say the Sanskrit for a tiny portion, then the Sanskrit for in a small degree.
\end{quote}

\subsection*{You'll want to know: \sk{अर्धम्}}
\textbf{\sk{अर्धम्}} — \emph{ardham} — \textquotedblleft{}a half\textquotedblright{}.

\begin{quote}
\sk{यदि} \sk{भवान्} \sk{अर्धं} \sk{खादति}, \sk{तर्हि} \sk{अर्धं} \sk{मह्यं} \sk{ददातु।}
\end{quote}

\begin{quote}
\emph{yadi bhavān ardhaṃ khādati, tarhi ardhaṃ mahyaṃ dadātu.}
\end{quote}

\begin{quote}
If you eat half, give half to me.
\end{quote}

A neuter noun.

\begin{cousinweb}
\textbf{\sk{अर्धम्}} (\emph{ardham}) MW gives \emph{ardha} \textquotedblleft{}half, halved\textquotedblright{} and, as a noun, \textquotedblleft{}the half\textquotedblright{}; MW gives no derivation, so its origin is uncertain (MW).
\end{cousinweb}

\begin{grammarlens}[title={What you've built}]
One more word for saying what happens if.
\end{grammarlens}

\subsection*{Your turn}
\begin{itemize}
\raggedright
  \item \emph{Say it:} \emph{ardham}
  \item \emph{Say it:} \emph{ardham}, once more
  \item \emph{Say it:} \emph{yadi bhavān ardhaṃ khādati, tarhi ardhaṃ mahyaṃ dadātu.} — If you eat half, give half to me.
  \item \emph{Recall:} say the Sanskrit for a tiny portion, then the Sanskrit for in a small degree, then say \emph{ardham} again
\end{itemize}

\begin{tcolorbox}[breakable,colback=teal!4,colframe=teal!35!black,title={Before you move on}]
What does \sk{अर्धम्} mean? (\textquotedblleft{}A half\textquotedblright{}.) Say it once more.
\end{tcolorbox}
\section[\texorpdfstring{\sk{न्यूनः} (nyūnaḥ)}{न्यूनः (nyūnaḥ)}]{\sk{न्यूनः} (nyūnaḥ) — less; falling short}

\label{lesson:SA-C378-nyunah}



\begin{quote}
Before the new one: say the Sanskrit for in a small degree, then the Sanskrit for a half.
\end{quote}

\subsection*{You'll want to know: \sk{न्यूनः}}
\textbf{\sk{न्यूनः}} — \emph{nyūnaḥ} — \textquotedblleft{}less; falling short\textquotedblright{}.

\begin{quote}
\sk{यदि} \sk{अन्नं} \sk{न्यूनम्}, \sk{तर्हि} \sk{अहं} \sk{न} \sk{खादिष्यामि।}
\end{quote}

\begin{quote}
\emph{yadi annaṃ nyūnam, tarhi ahaṃ na khādiṣyāmi.}
\end{quote}

\begin{quote}
If there is too little food, I will not eat.
\end{quote}

\begin{cousinweb}
\textbf{\sk{न्यूनः}} (\emph{nyūnaḥ}) is built from \emph{ni} and \emph{ūna}, \textquotedblleft{}wanting, deficient\textquotedblright{}; MW gives \emph{ny-ūna} \textquotedblleft{}less, diminished, defective, deficient\textquotedblright{} (MW).
\end{cousinweb}

\begin{grammarlens}[title={What you've built}]
One more word for saying what happens if.
\end{grammarlens}

\subsection*{Your turn}
\begin{itemize}
\raggedright
  \item \emph{Say it:} \emph{nyūnaḥ}
  \item \emph{Say it:} \emph{nyūnaḥ}, once more
  \item \emph{Say it:} \emph{yadi annaṃ nyūnam, tarhi ahaṃ na khādiṣyāmi.} — If there is too little food, I will not eat.
  \item \emph{Recall:} say the Sanskrit for in a small degree, then the Sanskrit for a half, then say \emph{nyūnaḥ} again
\end{itemize}

\begin{tcolorbox}[breakable,colback=teal!4,colframe=teal!35!black,title={Before you move on}]
What does \sk{न्यूनः} mean? (\textquotedblleft{}Less; falling short\textquotedblright{}.) Say it once more.
\end{tcolorbox}
\section[\texorpdfstring{\sk{ह्रसति} (hrasati)}{ह्रसति (hrasati)}]{\sk{ह्रसति} (hrasati) — dwindles; grows less}

\label{lesson:SA-C378-hrasati}



\begin{quote}
Before the new one: say the Sanskrit for a half, then the Sanskrit for falling short.
\end{quote}

\subsection*{You'll want to know: \sk{ह्रसति}}
\textbf{\sk{ह्रसति}} — \emph{hrasati} — \textquotedblleft{}dwindles; grows less\textquotedblright{}.

\begin{quote}
\sk{यदि} \sk{वर्षा} \sk{न} \sk{भवति}, \sk{तर्हि} \sk{नद्याः} \sk{जलं} \sk{ह्रसति।}
\end{quote}

\begin{quote}
\emph{yadi varṣā na bhavati, tarhi nadyāḥ jalaṃ hrasati.}
\end{quote}

\begin{quote}
If it does not rain, the river's water dwindles.
\end{quote}

Past (it …): \sk{जलं} \sk{ह्रसितवत्}, \emph{jalaṃ hrasitavat}.

\begin{cousinweb}
\textbf{\sk{ह्रसति}} (\emph{hrasati}) comes from the root √hras, \textquotedblleft{}to become short or small, to be diminished\textquotedblright{}; MW gives \emph{hrasita} \textquotedblleft{}shortened, curtailed, diminished\textquotedblright{} (MW).
\end{cousinweb}

\begin{grammarlens}[title={What you've built}]
One more word for saying what happens if.
\end{grammarlens}

\subsection*{Your turn}
\begin{itemize}
\raggedright
  \item \emph{Say it:} \emph{hrasati}
  \item \emph{Say it:} \emph{hrasati}, once more
  \item \emph{Say it:} \emph{yadi varṣā na bhavati, tarhi nadyāḥ jalaṃ hrasati.} — If it does not rain, the river's water dwindles.
  \item \emph{Recall:} say the Sanskrit for a half, then the Sanskrit for falling short, then say \emph{hrasati} again
\end{itemize}

\begin{tcolorbox}[breakable,colback=teal!4,colframe=teal!35!black,title={Before you move on}]
What does \sk{ह्रसति} mean? (\textquotedblleft{}Dwindles; grows less\textquotedblright{}.) Say it once more.
\end{tcolorbox}
